% ===== §11 Conclusions =====
\section{Conclusions: Coupling Degree as the True Measure of Luxury}
\label{sec:conclusions}

\noindent
This paper began with a question that seemed almost frivolous---which is more luxurious, a diamond ring or a flower on a path?---and has arrived at an answer that overturns the dominant framework of luxury theory and generates a justice theory of intimate gift-giving. In keeping with the series' commitment to the polyphonic conclusion form, this section presents what each of the paper's frameworks can and cannot conclude, ends with the structural closure of the three concentric circles, and explains why no synthetic conclusion is offered.

\subsection{What the Three Doctrines Establish}
\label{subsec:doctrines_establish}

\noindent
The three doctrines together establish that luxury in intimate relations is not a single phenomenon but a layered one, in which three distinct dimensions of value can be present in different combinations.

The first doctrine establishes that genuine luxury objects carry a double temporality (geological and artisanal) and achieve singularity through material non-replicability, biographical density, and relational contextualisation. The first doctrine cannot, however, explain why a flower---which has neither geological permanence nor artisanal craft---can be more luxurious than a diamond. For that, the third doctrine is required.

The second doctrine establishes that luxury in consumer society functions primarily through sign exchange value: as a differential position in the symbolic order, a phallic signifier enacting the logic of having and lacking, a socially legible signal of economic commitment. The second doctrine cannot, however, distinguish genuine from spurious intimate luxury, because it has eliminated the real register from its theoretical framework. For that distinction, the third doctrine is again required.

The third doctrine establishes that the true measure of luxury in intimate relations is coupling degree---the degree to which a subject's spatiotemporal structure is embedded in an object through existence-attestation. The third doctrine explains what the first two cannot: why the flower can outshine the diamond (high coupling degree at low economic value), why the proxy-purchased diamond is relationally impoverished (high economic value at low coupling degree), and why the long-saved ring is doubly luxurious (high coupling degree and high economic value reinforcing each other).

\subsection{What the Justice Theory Establishes}
\label{subsec:justice_establishes}

\noindent
The justice theory of choice establishes that the selection of objects in intimate relations is a value declaration with justice dimensions, and that justice in this domain concerns the integrity of existence-attestation rather than the distribution of goods.

It establishes that subject-embedding misrecognition---in its three forms of giver's self-misrecognition, receiver's induced misrecognition, and social collective misrecognition---is the central justice problem of relational luxury, and that the social collective misrecognition that equates economic value with coupling degree is the cultural condition that renders the flower invisible.

It establishes that equivalent agency---the substitution of one form of existence-attestation for another that the subject cannot provide---is structurally heterogeneous but may be just in the generative sense if four necessary conditions are met (symmetric generability, real constraint, genuine embedding, transparency), while insisting that these conditions are necessary but not sufficient, and that residual justice questions remain even when all four are satisfied.

It establishes that selection justice---the choice of coupling degree over sign exchange value as the primary criterion---is a justice stance: a refusal to use exchangeable symbols to bear the inechangeable relation, a reservation of space for the real register within the symbolic logic of consumer society.

The justice theory cannot, however, provide an algorithm for just choice. The assessment of coupling degree, the evaluation of the four conditions of equivalent agency, the judgement of whether genuine embedding has occurred---all of these require the cultivated sensibility that the theory describes but cannot mechanise. Justice in relational luxury is not a calculation but a practice, requiring the ongoing exercise of sensibility rather than the application of a rule.

\subsection{The Limits and Open Questions}
\label{subsec:limits}

\noindent
Several questions raised by the paper remain open.

\emb{The measurement of coupling degree}: The four dimensions of coupling degree have been characterised qualitatively, but the precise relationship between them---how they combine to constitute the overall coupling degree, whether they are commensurable, how trade-offs between them are to be assessed---remains underspecified. A more precise account of the structure of coupling degree is a task for future work.

\emb{The structure of existence-attestation}: The concept of existence-attestation has been developed through Ricoeur's account of attestation and the indexical analysis of \pinkref{Paper~XV} \citep{paperxv}, but the precise ontological mechanism by which a subject's spatiotemporal structure becomes embedded in an object---how the embedding occurs, what it consists in, how it persists---remains philosophically underdeveloped. This connects to the open questions about the Relational Isomorphism Principle raised in \pinkref{Paper~XV}: the embedding of subject in object is a case of the cross-register relation between the symbolic and the real that the series has not fully resolved.

\emb{The cultural specificity of the framework}: The framework has been developed primarily through Western and East Asian examples, and its applicability to other cultural traditions of luxury and gift remains to be examined. The cross-cultural analysis of \xref{sec:crosscultural} is a beginning, not a comprehensive account.

\emb{The application to non-material attestation}: The paper has focused on objects---material things in which a subject's existence can be embedded. But existence-attestation can occur through non-material media: through acts, words, presence, time given. The relationship between material relational luxury (the focus of this paper) and non-material existence-attestation (the act of presence, the gift of time) remains to be developed.

\subsection{Three Concentric Circles}
\label{subsec:circles}

\noindent
The paper closes, as the preceding papers in the series have closed, with the image of three concentric circles.

The \emb{innermost circle} is the specific object given in a specific relation: the flower on the path, the long-saved ring, the jade nourished by its wearer. This is the level of the concrete gift, the material thing in which a subject's existence is embedded and through which it is received. It is the level at which relational luxury is lived---in the giving and receiving of specific objects in specific relational moments.

The \emb{middle circle} is the political economy of luxury: the systems of production, distribution, and signification that determine what counts as luxury, who has access to it, and how it functions in the social symbolic order. This is the level of the three doctrines' analysis---the level at which the aesthetic, the symbolic, and the relational dimensions of luxury are produced and contested.

The \emb{outermost circle} is the ontological framework that grounds the entire analysis: the relational ontology in which the real is relational, the subject is co-constituted through the relational field, and existence-attestation is possible because a subject's spatiotemporal structure can be genuinely embedded in the objects they touch with care. This is the level at which the GRB series has been working throughout, and at which the concept of relational luxury finds its ultimate grounding.

The three circles are mutually constitutive: the specific gift is embedded in the political economy of luxury, and both are grounded in the ontological framework that makes existence-attestation possible. A complete account of relational luxury must attend to all three.

\subsection{Why No Synthetic Conclusion}
\label{subsec:no_synthesis}

\noindent
As in the preceding papers of the series, this paper offers no synthetic conclusion---no final paragraph that resolves all the analysis into a single unified statement. The reasons are by now familiar, but they take a specific form in the context of relational luxury.

A synthetic conclusion would imply that the theory of relational luxury is complete---that coupling degree has been fully specified, that the justice theory has been fully worked out, that the relationship between the three doctrines has been finally resolved. But the paper has been honest about what it does not know: the precise structure of coupling degree, the ontological mechanism of existence-attestation, the cultural variability of the framework, the relationship between material and non-material attestation. These are not loose ends but genuine open questions.

More fundamentally, a synthetic conclusion would betray the paper's central insight. Relational luxury is constituted by the irreducible particularity of the specific existence-attestation---by what cannot be subsumed under any general category or reduced to any universal formula. A paper that concluded with a universal formula for relational luxury would be performing, in its own form, the reduction of the particular to the general that the theory of relational luxury exists to resist. The refusal of synthetic closure is the paper's enactment of its own thesis: that what matters most---the genuine embedding of one existence in the life of another---cannot be captured in a formula but must be encountered in its particularity, recognised through sensibility, and received with the poetic openness that completes the coupling.

The flower on the path cannot be reduced to a concept. Neither can the conclusion of a paper about it.
